% Chapter 15 -- Adapting FLAIR to Your Work (Part II, Reference).  Sources
% (Law 1): os/milton/desktop_db.c and os/flair/finder_desktop.c (DESKTOP.DB:
% the desktop icons' places and each Finder window's place, size and state,
% written on an icon drag, a window move or close, Restart, Shut Down and
% Exit to InitechDOS, read at start-up); os/milton/kmain.c
% (finder_desk_load_positions; a Finder window reopens at its recorded
% place; flair_live_apple_refresh: the Apple menu re-read each time it is
% pulled down, apple_menu_build lists the rows of spec/apple_menu_items.h
% whose program is on the volume; the clock read from the RTC, MenuClock);
% os/milton/rtc.c rtc_set (TIME and DATE at the prompt set the computer's
% clock); os/flair/finder_menu.c (Preferences, the View by-items other than
% by Icons, and Sharing always dimmed); spec/assets/color_canon.json (the
% desktop colour is fixed).  Captures: screens/t_*.mouse.  Model: Macintosh
% User's Guide 1991 Ch. 11 "Adapting Your Macintosh to Your Own Use"
% (apple-macintosh.md Sec 8.2), cut to what FLAIR ships.
% GATES: test-desktop-db test-flair-desktop-icons test-flair-disk-windows test-flair-drag test-apple-menu test-flair-apple-menu test-flair-menubar-clock test-datetime test-finder-menu test-pubs-flair-screens
\Gates{test-desktop-db, test-flair-desktop-icons, test-flair-disk-windows,
  test-flair-drag, test-apple-menu, test-flair-apple-menu,
  test-flair-menubar-clock, test-datetime, test-finder-menu,
  test-pubs-flair-screens}
\chapter{Adapting FLAIR to Your Work}
\label{ch:adapting}

This chapter describes what you can arrange to suit the way you work: the
places of the icons and windows, the accessories in the Apple menu, and the
clock.  It also lists what this version of FLAIR does not let you change.

\section{How FLAIR remembers your arrangement}
\index{DESKTOP.DB}\index{arrangement of the desktop}
FLAIR keeps a record of your desktop in the file \file{DESKTOP.DB} on the
data disk: the place of the disk icon, and the place, size and state of
each Finder window.  It brings the record up to date when you:
\begin{itemize}
\item move an icon on the desktop;
\item move a Finder window or close one;
\item choose \menuitem{Restart}, \menuitem{Shut Down} or \menuitem{Exit to
  InitechDOS}.
\end{itemize}
When FLAIR starts, it reads the record and puts things back where you left
them.  The record belongs to the disk: the same data disk shows the same
arrangement on any computer.

If \file{DESKTOP.DB} is missing or cannot be read, FLAIR starts with the
icons in their usual places and makes a new record.

\section{Arranging the desktop}
\index{desktop!arranging}
\SideHead{The disk icon}Drag the \texttt{INITECH} icon anywhere on the
desktop (Figure~\ref{fig:icondrag}).  It stays there, and is there the next
time FLAIR starts.

\SideHead{The Trash}The Trash always stays at the lower right of the
screen.

\SideHead{HELLO and NOTES}The practice windows open in the same places each
time FLAIR starts.  You can move them while you work, but their places are
not recorded.

\section{Where windows open}
\index{windows!where they open}
A Finder window opens where you last left it.  The first time a folder is
opened, its window is placed a little lower and to the right of the last
one opened, so that their title bars can all be seen.  From then on, it
opens at the place and size it had when you moved it or closed it.
\begin{TaskSteps}
\Step Open the \texttt{INITECH} icon.
\Step Drag the \texttt{Drive A Files} window by its title bar down and to
the right.
\Step Click its close box.
\Step Double-click the \texttt{INITECH} icon again.
\Result{The window opens where you put it.}
\end{TaskSteps}

\FlairScreen[label=fig:twinplace]{t_winplace}{Drive A Files opened again
at the place it was moved to}

The icons inside a Finder window are not recorded.  Each time the window
opens, its icons are laid out in rows from the top left.

\section{The accessories in the Apple menu}
\index{Apple menu!accessories in}\index{accessories!adding and removing}
The Apple menu lists an accessory only while its program is in the
\file{APPS} folder of the data disk.  FLAIR looks at the disk each time you
pull the menu down, so a change shows at once.

\SideHead{To remove an accessory}Drag its program out of the \file{APPS}
folder, onto the \texttt{INITECH} icon or into another folder.  The
Calculator's program is \file{CALCULAT.EXE}.

\FlairScreen[scale=4, label=fig:tnoacc]{t_noacc}{The Apple menu after
CALCULAT.EXE was moved out of APPS}

\SideHead{To put it back}Drag its program back into the \file{APPS}
folder.  Its command returns to the Apple menu.

When no accessory is installed, the Apple menu holds only
\menuitem{About This Computer...}.

\section{Setting the clock}
\index{clock!setting}\index{TIME}\index{DATE}
The clock in the menu bar shows the time kept by the computer's own clock,
which runs even while the computer is turned off.  FLAIR has no command to
set it.  To set the clock, use the \file{TIME} command at the InitechDOS
prompt, and \file{DATE} to set the date.
\begin{TaskSteps}
\Step Choose \menuitem{Exit to InitechDOS} from the \menuname{Special} menu
and click \texttt{Exit}.
\Result{InitechDOS shows the date and the time and asks for new ones.}
\Step Type the correct date and press \key{Enter}, then type the correct
time and press \key{Enter}.
\Result{You can also set them later with \file{DATE} and \file{TIME},
followed by the new date or time.}
\end{TaskSteps}

\FlairText[fig:ttime]{t_time}{The computer's clock set to 1:45 pm with TIME}

\begin{TaskSteps}[resume]
\Step Restart the computer.
\Result{FLAIR starts, and its clock shows the time you set.}
\end{TaskSteps}

The clock shows hours and minutes, with \texttt{am} or \texttt{pm}.  It
cannot be changed to the 24-hour form, and it cannot be hidden.

\section{What you cannot change}
\index{Preferences command}
This version of FLAIR has no control panels.  The following are fixed:
\begin{itemize}
\item the colour of the desktop, the look of the windows and the
  typefaces of the menus;
\item the speed of the pointer and the double-click time;
\item the way Finder windows show their contents: as icons only.  In the
  \menuname{View} menu, \menuitem{by Icons} is the one view; the other
  views are dimmed;
\item the Finder's preferences: \menuitem{Preferences} in the
  \menuname{Edit} menu is dimmed.
\end{itemize}
The keys that FLAIR reserves are listed in Chapter~\ref{ch:keyboard}.
